\documentclass[aspectratio=169]{beamer}

% =========================
% Encoding, fonts, language
% =========================
\usepackage[T1]{fontenc}
\usepackage[utf8]{inputenc}
\usepackage[english,portuguese]{babel}
\selectlanguage{english}
\usepackage{lmodern}
\usepackage{microtype}

% =========================
% Graphics, math, tables
% =========================
\usepackage{graphicx}
\usepackage{amsmath,amssymb}
\usepackage{booktabs}
\usepackage{multirow}
\usepackage{xcolor,colortbl}
\usepackage{pgfplots}
\pgfplotsset{compat=1.17}
\usetikzlibrary{shapes.geometric,arrows,arrows.meta,positioning,fit,backgrounds,calc}

% =========================
% Theme & beamer housekeeping
% =========================
\usetheme{LACdefesa}
\setbeamertemplate{navigation symbols}{}

\setbeamerfont{title}{series=\bfseries,size=\Large}
\setbeamerfont{frametitle}{series=\bfseries}
\setbeamertemplate{itemize items}[circle]
\setbeamertemplate{enumerate items}[default]

% Logos del footer
\renewcommand{\LACLogoLeft}{images/logo-puc.png}
\renewcommand{\LACLogoCenter}{images/dept.png}

\titlegraphic{\includegraphics[height=1.0cm,keepaspectratio]{images/logo-puc.png}}

% =========================
% Colors & helpers
% =========================
\definecolor{PUCBlue}{HTML}{003C71}
\definecolor{Accent1}{HTML}{2A9D8F}
\definecolor{Accent2}{HTML}{E76F51}
\definecolor{Accent3}{HTML}{6A4C93}
\setbeamercolor{structure}{fg=PUCBlue}
\renewcommand{\arraystretch}{1.2}

% Macros
\newcommand{\good}[1]{\textcolor{Accent1}{\textbf{#1}}}
\newcommand{\warn}[1]{\textcolor{Accent2}{\textbf{#1}}}
\newcommand{\tbnote}[1]{\par{\tiny\color{gray}#1}}

% Colores del experimental design (de la proposta)
\definecolor{lowlabel}{RGB}{225, 245, 238}
\definecolor{highlabel}{RGB}{250, 236, 231}
\definecolor{sslpurple}{RGB}{238, 237, 254}
\definecolor{graybox}{RGB}{241, 239, 232}
\definecolor{accentgreen}{RGB}{29, 158, 117}
\definecolor{accentorange}{RGB}{216, 90, 48}
\definecolor{pool1}{RGB}{232, 225, 247}
\definecolor{pool2}{RGB}{210, 196, 240}
\definecolor{pool3}{RGB}{182, 160, 228}

% TikZ macro: experimental design figure (reutilizado de la proposta)
\newcommand{\makefigone}{%
\begin{tikzpicture}[
  x=1cm, y=1cm, font=\sffamily\footnotesize,
  every node/.style={inner sep=2.5pt},
  panel/.style={draw=gray!35, rounded corners=4pt, fill=gray!3},
  box/.style={draw, rounded corners=3pt, minimum width=3.1cm, minimum height=1.0cm, text width=2.8cm, align=center},
  sslbox/.style={draw, rounded corners=3pt, fill=sslpurple, draw=violet!60, minimum width=2.4cm, minimum height=0.55cm, align=center},
  >={Stealth[length=4pt]}
]
\draw[panel] (-8.0, 0.15) rectangle (8.0, -3.20);
\draw[panel] (-8.0, -3.55) rectangle (-0.2, -8.15);
\draw[panel] ( 0.2, -3.55) rectangle ( 8.0, -8.15);
\node[font=\sffamily\small\bfseries, anchor=west] at (-7.55, -0.28) {(A) Regime de quantidade de labels};
\node[draw=gray!60, rounded corners=2pt, fill=white, align=center, font=\scriptsize, text width=1.5cm, minimum height=1.0cm] (unmsub) at (-6.25, -1.25) {\textbf{Labels}\\\textbf{reduzidos}\\66, 111, 174 fr.};
\node[box, fill=lowlabel, draw=teal!60] (unm) at (-2.55, -1.25) {\textbf{UNM}\\[1pt] \scriptsize 218 frames de treino\\[-1pt] \scriptsize (23 pacientes)};
\node[box, fill=highlabel, draw=orange!60] (inca) at (2.55, -1.25) {\textbf{INCA}\\[1pt] \scriptsize 634 frames de treino\\[-1pt] \scriptsize (75 pacientes)};
\node[draw=gray!60, rounded corners=2pt, fill=white, align=center, font=\scriptsize, text width=1.5cm, minimum height=1.0cm] (incasub) at (6.25, -1.25) {\textbf{Labels}\\\textbf{reduzidos}\\56, 132, 283 fr.};
\draw[->, dashed, thin, gray!70!black] (unm.west) -- (unmsub.east);
\draw[->, dashed, thin, gray!70!black] (inca.east) -- (incasub.west);
\node[fill=white, inner sep=2pt, font=\scriptsize\itshape, text=gray!55!black] at ($(unm.west)!0.5!(unmsub.east) + (0,0.25)$) {subconjunto};
\node[fill=white, inner sep=2pt, font=\scriptsize\itshape, text=gray!55!black] at ($(inca.east)!0.5!(incasub.west) + (0,0.25)$) {subconjunto};
\node[sslbox] (sslq) at (0, -2.75) {\scriptsize SSL ajuda?};
\coordinate (joinL) at (-1.80, -2.75); \coordinate (joinR) at ( 1.80, -2.75);
\draw[thin] (unm.south) |- (joinL); \draw[thin] (inca.south) |- (joinR);
\draw[dashed, thin, gray!70!black] (unmsub.south) |- (joinL);
\draw[dashed, thin, gray!70!black] (incasub.south) |- (joinR);
\draw[->, thin] (joinL) -- (sslq.west); \draw[->, thin] (joinR) -- (sslq.east);
\node[font=\sffamily\small\bfseries, anchor=west] at (-7.80, -3.88) {(B) Tamanho do pool sem anotação};
\draw[dashed, gray!35, very thin] (-4.00, -4.50) -- (-4.00, -7.40);
\draw[thick, gray!55] (-6.60, -4.55) -- (-1.40, -4.55);
\draw[thick, gray!55] (-6.60, -4.48) -- (-6.60, -4.62);
\draw[thick, gray!55] (-1.40, -4.48) -- (-1.40, -4.62);
\node[font=\scriptsize, gray, anchor=west] at (-7.80, -4.20) {Linha do tempo do vídeo};
\fill[violet!80] (-4.07, -4.65) rectangle (-3.93, -4.45);
\draw[-{Stealth[length=3pt]}, gray!65, thin, shorten >=2pt] (-3.00, -4.82) -- (-3.85, -4.58);
\node[font=\scriptsize, gray!80!black, anchor=west] at (-2.95, -4.82) {Frame anotado};
\node[font=\scriptsize\bfseries, violet!45!black, anchor=west] at (-7.80, -5.20) {$r{=}3$};
\node[font=\scriptsize, gray, anchor=west] at (-7.00, -5.20) {$\pm 100$ ms};
\fill[pool1, draw=violet!45, thin] (-4.30, -5.32) rectangle (-3.70, -5.08);
\node[font=\scriptsize\bfseries, violet!55!black, anchor=west] at (-7.80, -5.90) {$r{=}10$};
\node[font=\scriptsize, gray, anchor=west] at (-7.00, -5.90) {$\pm 333$ ms};
\fill[pool2, draw=violet!55, thin] (-4.80, -6.02) rectangle (-3.20, -5.78);
\node[font=\scriptsize\bfseries, violet!70!black, anchor=west] at (-7.80, -6.60) {$r{=}20$};
\node[font=\scriptsize, gray, anchor=west] at (-7.00, -6.60) {$\pm 667$ ms};
\fill[pool3, draw=violet!65, thin] (-5.60, -6.72) rectangle (-2.40, -6.48);
\fill[graybox, draw=gray!55, thin] (-6.60, -7.42) rectangle (-1.40, -7.18);
\node[font=\scriptsize\itshape, gray!70!black] at (-4.00, -7.30) {todos os frames laterais (74{,}774)};
\draw[<->, gray!70, thin] (-1.05, -5.20) -- (-1.05, -7.30);
\node[font=\scriptsize, gray, rotate=90, align=center] at (-0.80, -6.25) {Tamanho do pool crescente};
\node[font=\sffamily\small\bfseries, anchor=west] at (0.55, -3.88) {(C) Política de seleção de frames};
\node[font=\scriptsize\bfseries, accentgreen, anchor=west] at (0.55, -4.40) {Proximidade temporal};
\node[font=\scriptsize, gray!75!black, anchor=west] at (0.55, -4.80) {Frames próximos ao anotado};
\draw[thick, gray!40] (0.55, -5.30) -- (7.60, -5.30);
\draw[thick, gray!55] (0.55, -5.22) -- (0.55, -5.38);
\draw[thick, gray!55] (7.60, -5.22) -- (7.60, -5.38);
\fill[violet!80] (4.02, -5.48) rectangle (4.18, -5.12);
\foreach \x in {3.5, 3.70, 3.90, 4.50, 4.3, 4.5, 4.7} {\fill[accentgreen] (\x-0.08, -5.48) rectangle (\x+0.08, -5.12);}
\node[font=\scriptsize\bfseries, accentorange, anchor=west] at (0.55, -6.00) {Aleatório, mesmo tamanho};
\node[font=\scriptsize, gray!75!black, anchor=west] at (0.55, -6.40) {Mesma quantidade, distribuídos pelo vídeo};
\draw[thick, gray!40] (0.55, -6.90) -- (7.60, -6.90);
\draw[thick, gray!55] (0.55, -6.82) -- (0.55, -6.98);
\draw[thick, gray!55] (7.60, -6.82) -- (7.60, -6.98);
\fill[violet!80] (4.02, -7.08) rectangle (4.18, -6.72);
\foreach \x in {0.70, 2.40, 2.60, 3.10, 4.40, 6.80, 7.40} {\fill[accentorange] (\x-0.08, -7.08) rectangle (\x+0.08, -6.72);}
\fill[violet!80] (1.00, -7.72) rectangle (1.20, -7.52); \node[font=\scriptsize, anchor=west] at (1.28, -7.62) {Anotado};
\fill[accentgreen] (3.30, -7.72) rectangle (3.50, -7.52); \node[font=\scriptsize, anchor=west] at (3.58, -7.62) {Temporal};
\fill[accentorange] (5.70, -7.72) rectangle (5.90, -7.52); \node[font=\scriptsize, anchor=west] at (5.98, -7.62) {Aleatório};
\end{tikzpicture}%
}

% Contador para preservar el total de slides al entrar a backup
\newcounter{finalframe}

% =========================
% Meta (titulo, autor, etc.)
% =========================
\title{Semi-supervised learning\\ for cervical vertebra segmentation\\ in videofluoroscopic swallowing studies}
\subtitle{M.Sc. Proposal Defense}
\author{Sebastian Quispe Arias}
\institute{Pontif\'{i}cia Universidade Cat\'{o}lica do Rio de Janeiro\\
\hspace{15}Departamento de Inform\'{a}tica}
\date{June 2026}


\begin{document}


%% ========================================
%% SLIDE 1: TITLE (layout custom)
%% ========================================
\begin{frame}[plain]
  % Logos en esquinas superiores (izquierda PUC, derecha DI)
  \begin{tikzpicture}[remember picture, overlay]
    \node[anchor=north west, inner sep=0pt] at ([xshift=0.5cm, yshift=-0.4cm]current page.north west)
      {\includegraphics[height=1.1cm,keepaspectratio]{images/logo-puc.png}};
    \node[anchor=north east, inner sep=0pt] at ([xshift=-0.5cm, yshift=-0.4cm]current page.north east)
      {\includegraphics[height=1.1cm,keepaspectratio]{images/dept.png}};
  \end{tikzpicture}

  \vspace{1.6cm}

  \begin{center}
    % T\'{i}tulo principal
    {\large\bfseries Semi-supervised learning\\[0.1cm]
     for cervical vertebra segmentation\\[0.1cm]
     in videofluoroscopic swallowing studies\par}

    \vspace{0.3cm}
    {\small\itshape Defesa de Mestrado}

    \vspace{0.8cm}

    % Autor + fecha
    {\normalsize Sebastian Quispe Arias}\\[0.2cm]
    {\footnotesize Setembro de 2026}
  \end{center}

  \vfill

  % Advisor centrado abajo
  \begin{center}
    \small
    \textbf{Orientador:} Prof. Paulo Ivson Netto Santos \\
    \textbf{Coorientador:} Luiz Fernando Trindade Santos
  \end{center}
\end{frame}


%% ========================================
%% DIVIDER 1: INTRODUÇÃO

%% ----- NUEVO: bottom line up front, antes de la motivacion -----
\begin{frame}{O que este trabalho pergunta, e o que encontrou}
  \vspace{0.3em}
  \textbf{RQ1. Em que regime de anotação o SSL ajuda?}
  \begin{itemize}\footnotesize
    \item Só onde o supervisionado ainda não saturou. Em UNM, $0{,}801 \rightarrow 0{,}860$. Em INCA com todas as etiquetas, nada.
  \end{itemize}

  \smallskip
  \textbf{RQ2. Quanto importa o pool de frames não rotulados?}
  \begin{itemize}\footnotesize
    \item Importa, e os dois métodos reagem em direções opostas.
  \end{itemize}

  \smallskip
  \textbf{RQ3. A proximidade temporal ajuda a escolher esses frames?}
  \begin{itemize}\footnotesize
    \item Não. Não se separa de uma escolha aleatória do mesmo tamanho.
  \end{itemize}

  \smallskip
  \textbf{E a jusante?}
  \begin{itemize}\footnotesize
    \item Uma máscara melhor coloca a região de interesse mais perto, mas não move o escalar C2--C4: o limite está na regra de pós-processamento, não na máscara.
  \end{itemize}

  \note{%
    Treinta segundos. Es el mapa para que sepan qu\'{e} van a juzgar.

    Para Geraldo, que no estuvo en la propuesta, este slide es lo que le permite colgar todo lo dem\'{a}s. Para Alberto y C\'{e}sar, es el recordatorio de d\'{o}nde qued\'{o} el trabajo.

    NO desarrollar aqu\'{i}. Cada l\'{i}nea tiene su slide de resultados m\'{a}s adelante.
  }
\end{frame}

%% ========================================
\begin{frame}[plain]
  \centering
  \vfill
  {\Huge\bfseries Introdução}
  \vfill
\end{frame}


%% ========================================
%% SLIDE 2: MOTIVATION
%% ========================================
\begin{frame}{Disfagia e VFSS}
  \begin{columns}[T,totalwidth=\textwidth]
    \column{0.55\textwidth}
      \textbf{Disfagia}
      \begin{itemize}
        \item Dificuldade para engolir
        \item Frequente em câncer de cabeça e pescoço
        \item Aspiração silenciosa $\rightarrow$ pneumonia, mortalidade
      \end{itemize}

      \bigskip

      \textbf{Videofluoroscopic swallowing study (VFSS)}
      \begin{itemize}
        \item Exame padrão para disfagia
        \item Vídeo de raios-X lateral durante a deglutição
        \item Vértebras cervicais C2--C4 como referência anatômica
      \end{itemize}

    \column{0.45\textwidth}
      \centering
      \vspace{0.6cm}
      \includegraphics[width=0.9\textwidth,height=5.5cm,keepaspectratio]{images/vfss_frame.png}
  \end{columns}

  \note{%
    Empezar: dysphagia es la dificultad para tragar. Muy com\'{u}n en pacientes con c\'{a}ncer de cabeza y cuello, donde radioterapia y cirug\'{i}a da\~{n}an las estructuras involucradas.

    Lo grave: aspiraci\'{o}n silenciosa. El paciente no tose pero el alimento o l\'{i}quido va al pulm\'{o}n. Causa neumon\'{i}a, desnutrici\'{o}n, y mortalidad.

    Necesitamos un examen objetivo: VFSS. V\'{i}deo de rayos X mientras el paciente traga un contraste (bario). Permite ver la deglutici\'{o}n din\'{a}micamente.

    Las v\'{e}rtebras cervicales C2--C4 son la referencia anat\'{o}mica clave para normalizar mediciones (escala C2-C4, tracking del hueso hioides, definir regiones de inter\'{e}s).

    Aqu\'{i} est\'{a} la imagen [se\~{n}alar]: se ve un frame lateral con las v\'{e}rtebras marcadas en verde.
  }
\end{frame}


%% ========================================
%% SLIDE 3: THE ANNOTATION BOTTLENECK
%% ========================================
\begin{frame}{O gargalo da anotação}
  \begin{columns}[T,totalwidth=\textwidth]
    \column{0.55\textwidth}
      \textbf{Anotação pixel a pixel é cara}
      \begin{itemize}
        \item Requer expertise clínica
        \item Vários minutos por frame
        \item Variabilidade entre anotadores (Dice $\approx$ 0,85)
        \item Restrições de privacidade limitam compartilhamento
      \end{itemize}

      \bigskip

      \textbf{Frames sem anotação são abundantes}
      \begin{itemize}
        \item Cada vídeo VFSS = milhares de frames
        \item Apenas alguns são anotados
        \item O resto é descartado no treinamento supervisionado
      \end{itemize}

    \column{0.45\textwidth}
      \centering
      \vspace{1.2cm}
      \begin{tikzpicture}[scale=0.8]
        % Eje
        \draw[->] (0,0) -- (5.2,0) node[right] {\scriptsize frames};
        \draw[->] (0,0) -- (0,4.2);

        % UNM
        \node[anchor=east] at (-0.1,3.5) {\scriptsize \textbf{UNM}};
        \fill[Accent2] (0,3.2) rectangle (0.025,3.8);
        \node[anchor=west, font=\tiny] at (0.05,3.5) {218 anotados};
        \fill[PUCBlue] (0,2.5) rectangle (5,3.1);
        \node[anchor=west, font=\tiny, white] at (0.1,2.8) {74{,}774 sem anotação};

        % INCA
        \node[anchor=east] at (-0.1,1.5) {\scriptsize \textbf{INCA}};
        \fill[Accent2] (0,1.2) rectangle (0.08,1.8);
        \node[anchor=west, font=\tiny] at (0.1,1.5) {634 anotados};
        \fill[PUCBlue] (0,0.5) rectangle (2.1,1.1);
        \node[anchor=west, font=\tiny, white] at (0.1,0.8) {31{,}260 sem anotação};

        % Leyenda
        \fill[Accent2] (0,-0.5) rectangle (0.3,-0.3);
        \node[anchor=west, font=\tiny] at (0.4,-0.4) {anotados};
        \fill[PUCBlue] (1.7,-0.5) rectangle (2,-0.3);
        \node[anchor=west, font=\tiny] at (2.1,-0.4) {sem anotação};
      \end{tikzpicture}
  \end{columns}

  \note{%
    El problema central: anotar VFSS al nivel de pixel es muy costoso. Requiere expertise cl\'{i}nica para distinguir los bordes de las v\'{e}rtebras entre estructuras superpuestas (mand\'{i}bula, hioides). Cada frame toma varios minutos.

    Adem\'{a}s, hay variabilidad entre anotadores: a\'{u}n expertos no se ponen 100\% de acuerdo. Medimos Dice de 0.854 en UNM y 0.806 en INCA.

    Y hay restricciones de privacidad: los datasets p\'{u}blicos son escasos.

    Pero hay algo importante: los frames sin anotar son abundantes. Cada v\'{i}deo VFSS tiene miles de frames, pero solo unos pocos se anotan. El resto se descarta en el entrenamiento supervisado tradicional.

    Aqu\'{i} est\'{a} la asimetr\'{i}a: en UNM tenemos 218 labeled vs 74,774 unlabeled. Ratio 343:1.

    Eso motiva nuestra pregunta: \textquestiondown podemos aprovechar esos frames no anotados?
  }
\end{frame}


%% ========================================
%% SLIDE 4: SSL IDEA
%% ========================================
\begin{frame}{Semi-supervised learning (SSL)}
  \begin{columns}[T,totalwidth=\textwidth]
    \column{0.5\textwidth}
      \textbf{SSL: combina frames anotados e sem anotação no treinamento}
      \begin{itemize}
        \item A cada iteração, combina perda supervisionada (com labels) $+$ perda não supervisionada (consistência ou pseudo-labeling)
        \item Funciona quando ambos os conjuntos compartilham a mesma distribuição (Oliver et al., 2018)
        \item Em VFSS: mesmos vídeos, mesma máquina, mesma anatomia \checkmark
      \end{itemize}

    \column{0.5\textwidth}
      \centering
      \vspace{1.0cm}
      \begin{tikzpicture}[
        scale=0.9, transform shape,
        labeledbox/.style={draw=Accent2!80, fill=Accent2!10, rounded corners=2pt, minimum width=1.4cm, minimum height=0.6cm, font=\scriptsize},
        unlabeledbox/.style={draw=PUCBlue!80, fill=PUCBlue!10, rounded corners=2pt, minimum width=1.4cm, minimum height=0.6cm, font=\scriptsize},
        model/.style={draw=black, fill=gray!15, rounded corners=3pt, minimum width=1.8cm, minimum height=1.2cm, font=\scriptsize, align=center}
      ]
        % Labeled (arriba izq)
        \node[labeledbox] (l1) at (0, 1.1) {anotados};
        \node[font=\tiny, gray, above=1pt of l1] {poucos};

        % Unlabeled stack (abajo izq) - dibujadas en orden para que u3 quede arriba (frontal)
        \node[unlabeledbox] (u3) at (0.2, -0.5) {};
        \node[unlabeledbox] (u2) at (0.1, -0.6) {};
        \node[unlabeledbox] (u1) at (0, -0.7) {sem anotação};
        \node[font=\tiny, gray, below=1pt of u1] {muitos};

        % Model (centro)
        \node[model] (m) at (3, 0.2) {Modelo de\\ segmentação};

        % Output (derecha) - propiedades inline
        \node[draw=Accent1, fill=Accent1!15, rounded corners=2pt, minimum width=1.7cm, minimum height=0.7cm, font=\scriptsize] (mask) at (5.5, 0.2) {máscara C2--C4};

        % Flechas - de cada caja al borde izquierdo del modelo, en el punto Y correspondiente
        \draw[->, thick] (l1.east) -- ([yshift=0.2cm]m.west);
        \draw[->, thick] (u1.east) -- ([yshift=-0.2cm]m.west);
        \draw[->, thick] (m.east) -- (mask.west);
      \end{tikzpicture}
  \end{columns}

  \note{%
    SSL es una clase de m\'{e}todos que aprovecha datos sin etiquetar junto con los etiquetados durante el entrenamiento.

    \textquestiondown Por qu\'{e} encaja bien con VFSS? Porque los frames sin etiquetar vienen de los MISMOS v\'{i}deos que los etiquetados (mismo dominio). Y los frames cercanos en tiempo suelen verse similares.

    En la literatura reciente de VFSS hay otras alternativas: active learning con expertos, foundation models tipo MedSAM, m\'{e}todos no supervisados.

    Pero el uso conjunto de frames etiquetados + no etiquetados del mismo v\'{i}deo (que es SSL cl\'{a}sico) no se ha estudiado sistem\'{a}ticamente para esta tarea.

    Por eso nuestro estudio se enfoca en SSL cl\'{a}sico: Pseudo-Labeling y Mean Teacher.
  }
\end{frame}


%% ========================================
%% SLIDE 3.5: SSL GAP IN VFSS (justificación)
%% ========================================
\begin{frame}{A lacuna de SSL em VFSS}
  \small
  \textbf{Busca no Scopus} (12/06/2026):
  \vspace{0.3em}

  \begin{block}{}
  \scriptsize\ttfamily
  (``semi-supervised'' OR ``pseudo-label*'' OR ``mean teacher'' OR \\
  \phantom{(}``consistency regularization'') \\
  AND \\
  (``VFSS'' OR ``videofluoroscop*'' OR ``swallowing study'')
  \end{block}

  \vspace{0.3em}
  $\rightarrow$ \textbf{19 documentos encontrados}
  \vspace{0.5em}

  Após revisão manual:
  \begin{itemize}
    \item 17 não são de VFSS (pâncreas, coxa, coluna lombar, ecocardiografia\ldots)
    \item 2 estudam mastigação/deglutição por áudio, não por imagem
    \item \textbf{0 abordam SSL para segmentação de vértebras em VFSS}
  \end{itemize}

  \vspace{0.4em}
  \textcolor{Accent1}{\textbf{Este trabalho preenche essa lacuna.}}
\end{frame}


%% ========================================
%% SLIDE 5: RESEARCH QUESTIONS
%% ========================================
\begin{frame}{Perguntas de pesquisa}
  \vspace{0.5em}

  \textbf{RQ1.} Em que regime de quantidade de labels o SSL melhora a segmentação?

  \vspace{0.8em}

  \textbf{RQ2.} Como o tamanho do pool de frames sem anotação afeta o desempenho do SSL?

  \vspace{0.8em}

  \textbf{RQ3.} A seleção por proximidade temporal supera a seleção aleatória do mesmo tamanho?

  \note{%
    Tres preguntas pr\'{a}cticas que no se han abordado conjuntamente en trabajos previos.

    RQ1: SSL es conocido que ayuda m\'{a}s cuando hay pocos labels. Pero el r\'{e}gimen exacto para VFSS no estaba caracterizado.

    RQ2: Tama\~{n}o del pool unlabeled importa. M\'{a}s no siempre es mejor (Fu et al. 2019 mostraron eso en laparoscop\'{i}a).

    RQ3: La estructura temporal de VFSS sugiere que frames vecinos al labeled podr\'{i}an ser \'{u}tiles. Comparamos contra selecci\'{o}n random del mismo tama\~{n}o.

    Para responder estas preguntas: dos m\'{e}todos SSL (Pseudo-Labeling y Mean Teacher), dos datasets (UNM y INCA).
  }
\end{frame}


%% ========================================
%% DIVIDER 2: SETUP EXPERIMENTAL
%% ========================================
\begin{frame}[plain]
  \centering
  \vfill
  {\Huge\bfseries Setup experimental}
  \vfill
\end{frame}


%% ========================================
%% SLIDE 6: DATASETS
%% ========================================
\begin{frame}{Datasets}
  \begin{columns}[T,totalwidth=\textwidth]
    \column{0.58\textwidth}
      \centering
      \footnotesize
      \begin{tabular}{@{}lcc@{}}
        \toprule
                                & \textbf{UNM}   & \textbf{INCA} \\
        \midrule
        Vídeos                  & 35             & 234 \\
        Pacientes               & 35             & 115 \\
        População               & Mista          & Câncer C\&P \\
        Frames anotados (treino)& 218            & 634 \\
        Frames sem anotação     & 74{,}774       & 31{,}260 \\
        Duração do vídeo (média)& 127\,s         & 6,6\,s \\
        Sem anotação : anotados & $\sim$343:1    & $\sim$49:1 \\
        \bottomrule
      \end{tabular}

    \column{0.42\textwidth}
      \textbf{Dois regimes complementares}
      \begin{itemize}
        \item UNM: poucos labels, pool sem anotação enorme
        \item INCA: mais labels, pool menor
      \end{itemize}

      \bigskip

      Os dois datasets usam o mesmo protocolo de anotação (ImageJ, supervisionado por especialista).

      \bigskip

      Dice com a especialista clínica:\\
      UNM 0,854 \quad INCA 0,806\\
      Entre anotadores (INCA, 898 frames): 0,874
  \end{columns}

  \note{%
    Dos datasets institucionales con estructuras distintas.

    UNM viene del TalkBank de la Universidad de New Mexico, p\'{u}blico. 35 v\'{i}deos largos (127 s en promedio) de pacientes con distintas condiciones de disfagia.

    INCA es del Instituto Nacional de C\^{a}ncer en Rio. 234 v\'{i}deos cortos (6.6 s cada uno), un solo trago por v\'{i}deo, pacientes con c\'{a}ncer de cabeza y cuello.

    Esta asimetr\'{i}a es importante: UNM ofrece pocos labels pero un pool unlabeled enorme (ratio 343:1). INCA ofrece m\'{a}s labels pero pool menor (49:1). Permite estudiar SSL en dos reg\'{i}menes diferentes.

    Para evaluar confiabilidad de las anotaciones, un especialista cl\'{i}nico independiente reanot\'{o} 20 frames de cada dataset. Acuerdo Dice de 0.854 en UNM y 0.806 en INCA.
  }
\end{frame}


%% ========================================
%% SLIDE 7: ARCHITECTURE AND BASELINE
%% ========================================
\begin{frame}{Arquitetura de segmentação}
  \begin{columns}[T,totalwidth=\textwidth]
    \column{0.5\textwidth}
      \textbf{U-Net++ com EfficientNet-B3}
      \begin{itemize}
        \item Encoder pré-treinado em ImageNet
        \item Backbone fixo em todos os experimentos SSL
        \item Entrada: $320 \times 320$ em escala de cinza
        \item Loss: BCE + Dice
      \end{itemize}

      \bigskip

      \textbf{Por que fixar o backbone?}
      \begin{itemize}
        \item Isolar o efeito do SSL da variabilidade da arquitetura
        \item Comparar contra 5 arquiteturas supervisionadas separadamente
      \end{itemize}

    \column{0.5\textwidth}
      \centering
      \vspace{0.8cm}
      \includegraphics[width=0.95\textwidth,height=5.5cm,keepaspectratio]{images/unetpp_arch.png}
      \tbnote{Adaptado de Zhou et al. (2018).}
  \end{columns}

  \note{%
    La arquitectura de segmentaci\'{o}n es U-Net++ con encoder EfficientNet-B3 preentrenado en ImageNet.

    Importante: la arquitectura es \textbf{fija} en todos los experimentos SSL. La raz\'{o}n es aislar el efecto del SSL del efecto de cambiar la arquitectura.

    Despu\'{e}s comparamos por separado contra 5 arquitecturas supervisadas (U-Net, DeepLabV3+, FPN, TransUNet, BiFPN-U-Net) y nnU-Net como referencia self-configuring.

    Frames en escala de grises a 320x320, normalizados a [0,1]. Loss combinado BCE + Dice.
  }
\end{frame}


%% ========================================
%% SLIDE 8: SSL METHODS
%% ========================================
\begin{frame}{Métodos avaliados: PL e MT}
  \begin{columns}[T,totalwidth=\textwidth]
    \column{0.5\textwidth}
      \textbf{Pseudo-Labeling (PL)}
      \begin{itemize}
        \item Teacher gera pseudo-labels duros onde está confiante: $p \geq 0{,}95$ ou $p \leq 0{,}05$
        \item Student aprende apenas dos pixels confiáveis
      \end{itemize}

      \bigskip

      \textbf{Mean Teacher (MT)}
      \begin{itemize}
        \item O mapa de probabilidades suave do teacher é o alvo de consistência
        \item Todos os pixels contribuem, incluindo os incertos
      \end{itemize}

    \column{0.5\textwidth}
      \centering
      \vspace{1.2cm}
      \begin{tikzpicture}[
        scale=0.85, transform shape,
        net/.style={draw=black, fill=PUCBlue!15, rounded corners=2pt, minimum width=1.6cm, minimum height=0.7cm, font=\scriptsize},
        aug/.style={font=\tiny, gray, fill=white, inner sep=1pt},
        out/.style={draw=Accent1!80, fill=Accent1!15, rounded corners=2pt, minimum width=1.5cm, minimum height=0.55cm, font=\tiny}
      ]
        % Input (izquierda centro)
        \node[font=\scriptsize] (inp) at (0, 0) {sem anotação};

        % Teacher (arriba)
        \node[net] (teacher) at (2.6, 1) {Teacher};
        \node[draw=Accent1, fill=Accent1!15, rounded corners=2pt, minimum width=1.6cm, minimum height=0.6cm, font=\tiny] (tout) at (5.2, 1) {previsão};

        % Student (abajo)
        \node[net] (student) at (2.6, -1) {Student};
        \node[draw=Accent1, fill=Accent1!15, rounded corners=2pt, minimum width=1.6cm, minimum height=0.6cm, font=\tiny] (sout) at (5.2, -1) {previsão};

        % Flechas input → nets
        \draw[->] (inp.east) -- node[aug, pos=0.55]{aug fraco} (teacher.west);
        \draw[->] (inp.east) -- node[aug, pos=0.55]{aug forte} (student.west);

        % Flechas nets → outs
        \draw[->] (teacher.east) -- (tout.west);
        \draw[->] (student.east) -- (sout.west);

        % EMA (student → teacher)
        \draw[->, dashed, gray] (student.north) to[bend right=40] node[font=\tiny, midway, right, gray]{EMA} (teacher.south);

        % Match loss
        \node[font=\tiny, align=center, gray] (loss) at (7.2, 0) {match\\(loss)};
        \draw[<->, dotted, gray!70] (tout.south) -- (loss.north);
        \draw[<->, dotted, gray!70] (sout.north) -- (loss.south);
      \end{tikzpicture}
  \end{columns}

  \note{%
    Dos m\'{e}todos representativos de familias distintas de SSL para segmentaci\'{o}n.

    Ambos comparten configuraci\'{o}n teacher-student. El teacher es un EMA (exponential moving average) de los pesos del student, con coeficiente 0.99.

    PL: el teacher genera pseudo-labels duras (binarizadas). Solo se usan los p\'{i}xeles donde el teacher est\'{a} muy confidente (probabilidad $\geq$ 0.95 o $\leq$ 0.05). Sigue FixMatch.

    MT: en lugar de binarizar, se usa la probabilidad continua del teacher como target. El student se entrena para coincidir con MSE. Todos los p\'{i}xeles contribuyen.

    \textquestiondown Cu\'{a}l es mejor? Depende del r\'{e}gimen, lo veremos en resultados.
  }
\end{frame}


%% ========================================
%% SLIDE 9: UNLABELED POOLS
%% ========================================
\begin{frame}{Desenho experimental}
  \centering
  \resizebox{!}{0.78\textheight}{\makefigone}

  \note{%
    Este slide es la figura completa del experimental design del paper.

    Panel A: las dos datasets (UNM, INCA) con sus reducciones de labels para los experimentos de label-efficiency.

    Panel B: construcci\'{o}n de pools unlabeled por radio temporal. Para cada frame etiquetado en el tiempo $t$, tomamos los frames cercanos dentro de una ventana $[t-r, t+r]$. Probamos seis radios: 3, 5, 7, 10, 15, 20. Y un pool grande con todos los frames laterales del video.

    Los radios cubren las escalas de tiempo relevantes de la degluci\'{o}n: r=3 son $\pm$100ms (un sub-evento), r=20 son $\pm$667ms (casi una transici\'{o}n far\'{i}ngea completa).

    Panel C: las dos pol\'{i}ticas de selecci\'{o}n. Temporal-neighbor (verde) selecciona frames cercanos al labeled. Size-matched random (naranja) toma el mismo n\'{u}mero pero distribuidos aleatoriamente en el video. Esto permite separar el efecto del tama\~{n}o del pool del efecto de la pol\'{i}tica.
  }
\end{frame}


%% ========================================
%% DIVIDER 3: RESULTADOS
%% ========================================
\begin{frame}[plain]
  \centering
  \vfill
  {\Huge\bfseries Resultados}
  \vfill
\end{frame}


%% ========================================
%% SLIDE 10: RESULT 1 - SSL ACROSS LABELS
%% ========================================
\begin{frame}{Resultado 1: quando o SSL ajuda}
  \begin{columns}[T,totalwidth=\textwidth]
    \column{0.55\textwidth}
      \centering
      \includegraphics[height=6.5cm,keepaspectratio]{images/fig3_label_efficiency_vertical.pdf}

    \column{0.45\textwidth}
      \textbf{Quando SSL ajuda}
      \begin{itemize}
        \item UNM (218): \phantom{0,000} $0{,}801 \to$ \good{$0{,}860$} \footnotesize ($\Delta = +0{,}059$)
        \item INCA 10\% (56): $0{,}844 \to$ \good{$0{,}865$} \footnotesize ($\Delta = +0{,}021$)
      \end{itemize}

      \bigskip

      \textbf{Quando SSL não ajuda}
      \begin{itemize}
        \item INCA full (634): $0{,}906 = 0{,}906$
      \end{itemize}

      \bigskip

      \textbf{MT preserva qualidade sob redução de labels}\\
      \footnotesize $-4\%$ F1 em INCA 10\%, $-1\%$ em UNM 25\%
  \end{columns}

  \note{%
    El primer resultado responde RQ1: SSL ayuda cuando la baseline supervisada no ha saturado.

    En UNM, con 218 frames etiquetados, la supervisada llega a F1 = 0.801. Mean Teacher con el pool all-lateral sube a 0.860, una mejora de +0.059 (significativa, Wilcoxon a nivel video).

    En INCA con el dataset completo (634 frames), no hay mejora: la supervisada ya est\'{a} en 0.906 y SSL la iguala. Ya satur\'{o}.

    Para verificar que no es una propiedad del dataset sino del r\'{e}gimen de etiquetas, redujimos INCA a 10\%, 25\%, 50\%. Solo a 10\% (56 frames) SSL vuelve a mejorar: MT llega a 0.865 vs 0.844 supervisado.

    Conclusi\'{o}n: SSL no es universal. Funciona cuando la supervisada est\'{a} bajo aproximadamente F1 = 0.85.
  }
\end{frame}


%% ========================================
%% SLIDE 11: RESULT 2 - POOL COMPOSITION
%% ========================================
\begin{frame}{Resultado 2: o tamanho do pool (RQ2)}
  \begin{columns}[T,totalwidth=\textwidth]
    \column{0.55\textwidth}
      \centering
      \includegraphics[height=6.5cm,keepaspectratio]{images/fig2_pool_size_curves_vertical.pdf}

    \column{0.45\textwidth}
      \textbf{Os dois métodos reagem ao contrário}
      \begin{itemize}
        \item PL: melhor com o pool \emph{pequeno} ($0,852$ em $r{=}3$), pior com o maior ($0,803$ em $r{=}20$)
        \item MT: melhor com o pool \emph{completo} ($0,860$), com $74{.}774$ frames
      \end{itemize}

      \bigskip

      \textbf{Por quê}
      \begin{itemize}
        \item PL compromete um rótulo duro por pixel: um erro confiante do professor reforça o erro
        \item MT segue um alvo contínuo, mais tolerante ao ruído
      \end{itemize}

      \bigskip
      \footnotesize\textbf{Usar todos os frames disponíveis não é a escolha segura.}
  \end{columns}

  \note{%
    El segundo resultado responde RQ2 y RQ3.

    RQ2 (tama\~{n}o del pool): los dos m\'{e}todos reaccionan distinto. PL prefiere pools peque\~{n}os: mejor en r=3 (1,257 frames). MT prefiere pools grandes: mejor en all-lateral (74,774 frames).

    \textquestiondown Por qu\'{e}? PL compromete una etiqueta dura por p\'{i}xel, entonces los errores del teacher son m\'{a}s da\~{n}inos en pools amplios. MT usa target continuo, es m\'{a}s robusto al ruido.

    RQ3 (pol\'{i}tica de selecci\'{o}n): la proximidad temporal no aporta ventaja significativa sobre selecci\'{o}n random del mismo tama\~{n}o. Tests Wilcoxon con p mayor a 0.05 en todas las configuraciones.

    Esto es contraintuitivo: la estructura temporal de VFSS pareci\'{a} sugerir que los vecinos temporales ayudar\'{i}an m\'{a}s. No fue el caso.
  }
\end{frame}

%% ----- NUEVO: RQ3 como resultado propio -----
\begin{frame}{Resultado 3: a proximidade temporal não ajuda (RQ3)}
  \vspace{0.3em}
  \textbf{Em cada tamanho de pool, a seleção temporal foi comparada com um controle aleatório do mesmo tamanho}

  \bigskip
  \centering
  \footnotesize
  \begin{tabular}{@{}lcccc@{}}
    \toprule
    \textbf{UNM, pseudo-rotulagem} & $r{=}3$ & $r{=}7$ & $r{=}10$ & $r{=}15$ \\
    \midrule
    Temporal   & $0,852$ & $0,830$ & $0,832$ & $0,840$ \\
    Aleatório  & $0,835$ & $0,819$ & $0,845$ & $0,846$ \\
    \bottomrule
  \end{tabular}

  \bigskip
  \raggedright
  \begin{itemize}\footnotesize
    \item O sinal muda com o raio: temporal ganha nos pequenos, perde nos grandes
    \item No raio de referência $r{=}10$, testes de Wilcoxon por paciente não detectaram diferença para nenhum dos dois métodos, em nenhuma semente
    \item Em INCA, as diferenças ficam abaixo de $0{,}001$ na maioria das configurações
  \end{itemize}

  \medskip
  \footnotesize\textbf{A estrutura temporal do vídeo não é, por si só, um critério útil para escolher frames.}

  \note{%
    Esta era la tercera pregunta de investigaci\'{o}n y la respuesta es NO. Es un resultado, no una nota al pie.

    La intuici\'{o}n dec\'{i}a lo contrario: en v\'{i}deo, los vecinos temporales deber\'{i}an ser los frames m\'{a}s \'{u}tiles. No lo son.

    Lo importante: no es que salga poco. Es que el SIGNO cambia con el radio. Gana en r=3 y r=7, pierde en r=10 y r=15. Eso es ruido, no un efecto.

    Y en r=10, que es el radio de referencia del trabajo, el Wilcoxon por paciente no detecta diferencia en ninguna semilla ni con ninguno de los dos m\'{e}todos.

    Si preguntan por qu\'{e} lo esperaba: porque los frames vecinos comparten anatom\'{i}a y el modelo ya la ve. Aportan poca informaci\'{o}n nueva.
  }
\end{frame}



%% ========================================
%% SLIDE 12: RESULT 3 - VS ARCHITECTURES
%% ========================================
\begin{frame}{Resultado 4: SSL vs arquiteturas supervisionadas}
  \begin{columns}[T,totalwidth=\textwidth]
    \column{0.55\textwidth}
      \centering
      \footnotesize
      \textbf{UNM, regime de labels completo (218)}
      \begin{tabular}{@{}lc@{}}
        \toprule
        \textbf{Método} & \textbf{F1} \\
        \midrule
        \multicolumn{2}{@{}l}{\emph{Supervisionado ($320{\times}320$)}} \\
        \quad BiFPN-U-Net(T)                & 0,759 \\
        \quad DeepLabV3+                    & 0,776 \\
        \quad FPN                           & 0,799 \\
        \quad U-Net++ (baseline)            & 0,801 \\
        \quad U-Net                         & 0,824 \\
        \quad TransUNet                     & 0,851 \\
        \midrule
        \multicolumn{2}{@{}l}{\emph{SSL ($320{\times}320$)}} \\
        \quad PL ($r{=}3$, U-Net++)         & \good{0,852} \\
        \quad MT (pool completo, U-Net++)   & \good{0,860} \\
        \midrule
        \multicolumn{2}{@{}l}{\emph{Auto-configurado ($1024{\times}1024$)}} \\
        \quad nnU-Net                       & 0,907 \\
        \bottomrule
      \end{tabular}

    \column{0.45\textwidth}
      SSL com U-Net++ alcança F1 comparável à melhor arquitetura supervisionada do pipeline (TransUNet, 0,851).

      \bigskip

      nnU-Net chega a F1 = 0,907, mas usa:
      \begin{itemize}
        \item Resolução diferente ($1024 \times 1024$)
        \item Treino mais longo, augmentation diferente
      \end{itemize}
  \end{columns}

  \note{%
    Tercer resultado: comparaci\'{o}n contra arquitecturas supervisadas.

    En UNM con 218 etiquetados, evaluamos 5 arquitecturas supervisadas. El rango va de BiFPN-U-Net(T) en 0.759 (sin pretraining) a TransUNet en 0.851, la m\'{a}s fuerte de las testeadas.

    SSL con U-Net++ alcanza 0.852 (PL) y 0.860 (MT). Iguala o supera la arquitectura supervisada m\'{a}s fuerte.

    Por separado, nnU-Net llega a 0.907. Pero hay que interpretarlo con cuidado: nnU-Net usa resoluci\'{o}n 1024x1024, entrenamiento mucho m\'{a}s largo, augmentation distinto. Es una comparaci\'{o}n a nivel de pipeline, no de paradigma de aprendizaje.

    El mensaje principal: SSL y la optimizaci\'{o}n del pipeline son direcciones complementarias. SSL mejora sin cambiar el backbone, nnU-Net adapta el pipeline al dataset.
  }
\end{frame}

%% ----- NUEVO: SSL sobre los otros backbones -----
\begin{frame}{SSL em outros backbones}
  \centering
  \footnotesize
  \begin{tabular}{@{}lccc@{}}
    \toprule
    \textbf{Backbone} & \textbf{Supervisionado} & \textbf{Mean Teacher} & \textbf{$\Delta$} \\
    \midrule
    U-Net             & $0,824$ & $0,851$ & $+0,027$ \\
    DeepLabV3+        & $0,776$ & $0,800$ & $+0,024$ \\
    U-Net++           & $0,801$ & $0,830$ & $+0,029$ \\
    FPN               & $0,799$ & $0,811$ & $+0,012$ \\
    BiFPN-U-Net(T)    & $0,759$ & $0,757$ & $-0,002$ \\
    TransUNet         & $0,851$ & $0,831$ & $-0,020$ \\
    \bottomrule
  \end{tabular}

  \vspace{0.8em}
  \footnotesize
  \textit{Mesma configuração para os seis: Mean Teacher, pool aleatório $r{=}15$.}

  \begin{itemize}\footnotesize
    \item Quatro dos seis ganham entre $0,012$ e $0,029$
    \item TransUNet, a melhor supervisionada, é a única que perde
  \end{itemize}

  \note{%
    Este frame responde a Paulo, que pregunt\'{o} d\'{o}nde estaba el SSL sobre todos los backbones.

    El punto: el uso de frames no etiquetados no est\'{a} atado a U-Net++. Cuatro de seis mejoran.

    Pero NO llega a todos. BiFPN no se mueve y TransUNet pierde 0,020. Y eso contradice la intuici\'{o}n de Jean en la propuesta, que esperaba que TransUNet, por ser m\'{a}s profunda, se beneficiara m\'{a}s.

    Si preguntan por qu\'{e}: no lo sabemos. Lo medimos y sale as\'{i}. Una hip\'{o}tesis es la saturaci\'{o}n: TransUNet ya parte de 0,851 supervisada, que es donde el SSL deja de ayudar en el resto del trabajo.

    OJO: el bloque supervisado es de abril y dos filas del bloque SSL son de septiembre. La conclusi\'{o}n aguanta por los dos caminos, est\'{a} verificado.
  }
\end{frame}

%% ----- NUEVO: la descomposicion de nnU-Net -----
\begin{frame}{nnU-Net: o que explica a diferença?}
  \begin{columns}[T,totalwidth=\textwidth]
    \column{0.46\textwidth}
      \textbf{nnU-Net como referência}
      \begin{itemize}\footnotesize
        \item $0,896$ (25\%) a $0,907$ (100\%)
        \item Acima de tudo o que treinamos
        \item Mas o pipeline difere em tudo: resolução, augmentation, duração, normalização
      \end{itemize}

      \bigskip
      \textbf{Cada diferença, uma por uma}
      \begin{itemize}\footnotesize
        \item Transplantada sozinha no nosso pipeline
        \item Medida contra o seu próprio controle, três sementes
      \end{itemize}

    \column{0.54\textwidth}
      \centering
      \footnotesize
      \begin{tabular}{@{}lcc@{}}
        \toprule
        \textbf{Configuração} & \textbf{$\Delta$ F1} & \textbf{Sementes} \\
        \midrule
        Resolução $1024^2$        & $+0,001 \pm 0,036$ & 2 de 3 \\
        Augmentation moderada     & $-0,003 \pm 0,038$ & 2 de 3 \\
        Augmentation completa     & $-0,017 \pm 0,014$ & 0 de 3 \\
        Normalização z-score      & $-0,012 \pm 0,064$ & 1 de 3 \\
        \bottomrule
      \end{tabular}

      \vspace{0.8em}
      \footnotesize
      \textbf{Nenhuma explica a diferença sozinha.}
      A única mudança consistente é uma \emph{queda} com a augmentation completa.
  \end{columns}

  \note{%
    Este frame responde a Alberto y a C\'{e}sar, que en la propuesta dijeron que la comparaci\'{o}n con nnU-Net no era justa.

    La respuesta no es defenderse: es descomponerla. Cogimos cada diferencia de configuraci\'{o}n y la trasplantamos sola a nuestro pipeline.

    Ninguna explica la diferencia. Y la \'{u}nica consistente, con las tres semillas en la misma direcci\'{o}n, es que la augmentation completa de nnU-Net EMPEORA.

    Los dos perfiles comparten toda la fotometr\'{i}a y solo difieren en la geometr\'{i}a, as\'{i} que la resta atribuye el da\~{n}o a la geometr\'{i}a: rotar 180 grados y espejar produce im\'{a}genes que el modelo nunca ve en test.

    Si preguntan por 480 de resoluci\'{o}n: no la probamos. Probamos la alternativa que estaba en cuesti\'{o}n, la de nnU-Net.
  }
\end{frame}



%% ========================================
%% SLIDE 13: QUALITATIVE RESULTS
%% ========================================
\begin{frame}{Resultados qualitativos}
  \centering
  \includegraphics[width=\textwidth,height=6.9cm,keepaspectratio]{images/fig5_qualitative_horizontal_annotated.pdf}

  \vspace{0.4em}
  \footnotesize
  \textcolor{Accent1}{Verde} = TP, \textcolor{Accent2}{vermelho} = FN, \textcolor{Accent3}{azul} = FP.

  \note{%
    Este slide es principalmente visual.

    La figura muestra ejemplos cualitativos en los tres reg\'{i}menes: UNM, INCA 10\% (pocos labels), INCA full.

    Filas: input frame, ground truth, predicci\'{o}n supervisada, predicci\'{o}n con Mean Teacher.

    Lo importante es que en los reg\'{i}menes donde SSL mejora cuantitativamente (UNM y INCA 10\%), tambi\'{e}n se ve cualitativamente: bordes m\'{a}s precisos, menos false negatives (zonas rojas) que el supervisado.

    En INCA full, donde SSL no mejora, las predicciones son similares.

    Este es un buen punto para que la banca observe los datos directamente. Yo destacar\'{i}a las mejoras visibles en UNM (primera columna).
  }
\end{frame}

%% ----- NUEVO: incerteza -----
\begin{frame}{Incerteza preditiva}
  \begin{columns}[T,totalwidth=\textwidth]
    \column{0.52\textwidth}
      \centering
      \includegraphics[width=\linewidth,height=5.2cm,keepaspectratio]{images/fig_uncertainty_vs_pool.pdf}

    \column{0.48\textwidth}
      \textbf{Onde está a incerteza}
      \begin{itemize}\footnotesize
        \item Entropia binária por pixel, numa banda de $5$~px em volta do contorno
        \item Banda $0,049$ \quad dentro das vértebras $0,029$
        \item Vãos entre vértebras $0,022$ \quad fundo $0,0005$
      \end{itemize}

      \bigskip
      \textbf{O SSL não deixa o modelo mais certo}
      \begin{itemize}\footnotesize
        \item $0,049 \rightarrow 0,048$; o que muda é \emph{onde} a incerteza fica
        \item O tamanho do pool também não muda: $0,041$ a $0,049$ de $1{.}257$ a $74{.}774$ frames
      \end{itemize}
  \end{columns}

  \note{%
    Este frame responde a Paulo, que dijo que no hab\'{i}a explicado c\'{o}mo se calcula la incertidumbre y pidi\'{o} un gr\'{a}fico contra el tama\~{n}o del pool.

    La medida: entrop\'{i}a binaria de la probabilidad predicha, por p\'{i}xel. M\'{a}ximo ln 2 = 0,693. O sea que 0,049 es el 7 por ciento del m\'{a}ximo posible.

    D\'{o}nde est\'{a}: pegada al contorno. El fondo pr\'{a}cticamente no duda.

    Y el resultado negativo, que es el interesante: el SSL NO baja la incertidumbre total. La MUEVE, de dentro de las v\'{e}rtebras hacia los vanos entre ellas. Y correlaciona mejor con el error, de -0,25 a -0,51.

    El gr\'{a}fico: el pool crece sesenta veces y la entrop\'{i}a no se mueve. Cambiar la semilla la mueve m\'{a}s que multiplicar el pool por sesenta.

    NO decir que la incertidumbre sirve para seleccionar frames. Lo probamos y el contraste preregistrado sale nulo.
  }
\end{frame}

%% ----- NUEVO: consecuencias practicas -----
\begin{frame}{Consequências práticas}
  \begin{columns}[T,totalwidth=\textwidth]
    \column{0.5\textwidth}
      \centering
      \includegraphics[width=\linewidth,height=4.6cm,keepaspectratio]{images/fig_roi_lee.pdf}

      \vspace{0.3em}
      {\scriptsize\itshape Um frame do test set. Os valores ao lado são a média sobre os 63 frames.}

    \column{0.5\textwidth}
      \textbf{Região de interesse dinâmica}
      \begin{itemize}\footnotesize
        \item A máscara só dá uma \emph{posição}: o centroide ancora a janela
        \item Supervisionado $6{,}09$~px do centroide de referência; Mean Teacher $4{,}62$~px
        \item MT mais perto em $44$ dos $63$ frames ($p<0{,}001$)
      \end{itemize}

      \bigskip
      \textbf{Escalar anatômico C2--C4}
      \begin{itemize}\footnotesize
        \item Erro mediano $9{,}3 \rightarrow 7{,}9$~px
        \item Mas a regra de pós-processamento tem um piso de $6{,}9$~px, medido sobre as máscaras de referência
        \item Máscaras melhores não movem mais o escalar
      \end{itemize}
  \end{columns}

  \note{%
    Este frame responde a C\'{e}sar, que pregunt\'{o} para qu\'{e} sirve cl\'{i}nicamente y si hace falta una segmentaci\'{o}n tan buena.

    Dos tareas aguas abajo, medidas.

    La ROI de Lee: la ventana tiene tama\~{n}o fijo, as\'{i} que la m\'{a}scara solo aporta una posici\'{o}n. Mean Teacher la coloca m\'{a}s cerca, y gana en 44 de 63 frames.

    El escalar C2-C4: aqu\'{i} la respuesta es incomoda y la doy yo antes de que la pregunten. La mejora en C4 NO se traslada al escalar. Solo baja un 15 por ciento.

    Por qu\'{e}: la regla que localiza las esquinas tiene un suelo de 6,9 p\'{i}xeles, medido aplic\'{a}ndola a las m\'{a}scaras de referencia, donde la segmentaci\'{o}n es correcta por construcci\'{o}n. Ese suelo se come el 87 por ciento del error.

    Conclusi\'{o}n honesta: para ese escalar, m\'{a}s ganancia tiene que venir de la regla o de predecir las esquinas directamente, no de mejores m\'{a}scaras.
  }
\end{frame}



%% ========================================
%% SLIDE 14: LIMITATIONS
%% ========================================

%% ----- NUEVO: contribuciones -----
\begin{frame}{Contribuições}
  \vspace{0.4em}
  \begin{itemize}
    \item Primeira avaliação sistemática de SSL clássico para segmentação das vértebras cervicais em VFSS, em dois datasets institucionais
    \medskip
    \item Um estudo controlado da \textbf{construção do pool não rotulado}, tamanho e política, com o resto do pipeline fixo, de modo que as diferenças sejam atribuíveis
    \medskip
    \item Comparação contra \textbf{seis arquiteturas supervisionadas} e contra nnU-Net, com a diferença decomposta configuração por configuração
    \medskip
    \item Duas tarefas \textbf{a jusante} medidas, e não apenas o Dice: o escalar anatômico C2--C4 e o posicionamento da região de interesse
  \end{itemize}

  \note{%
    Este slide no estaba y hace falta, sobre todo para Geraldo.

    Lo que hay que transmitir: esto NO es un algoritmo nuevo de SSL. Es un estudio emp\'{i}rico de c\'{o}mo se construye el pool, con el pipeline fijo para que las diferencias sean atribuibles.

    La cuarta es la que distingue el trabajo de un benchmark: no nos quedamos en el Dice, medimos si la m\'{a}scara mejor sirve para lo que el cl\'{i}nico hace despu\'{e}s.

    NO decir "primer trabajo que aplica SSL a VFSS". Decir "primera evaluaci\'{o}n sistem\'{a}tica", que es lo que est\'{a} sostenido.
  }
\end{frame}

\begin{frame}{Limitações}
  \begin{itemize}
    \item Backbone único (U-Net++ a $320 \times 320$); efeitos podem diferir com outros backbones
    \medskip
    \item FixMatch e cross-pseudo supervision não testadas, embora sejam refinamentos das mesmas duas famílias
    \medskip
    \item Test set pequeno em UNM (7 vídeos); INCA não é publicamente compartilhável
  \end{itemize}

  \note{%
    Tres limitaciones honestas.

    Primera: usamos un solo backbone fijo a 320x320. No sabemos si los efectos del tama\~{n}o del pool y selecci\'{o}n se mantienen con otros backbones.

    Segunda: m\'{e}todos m\'{a}s recientes (FixMatch, cross-pseudo supervision) y otras familias (adversariales, contrastive) no fueron testados.

    Tercera: el test set de UNM tiene solo 7 videos, lo que limita el poder estad\'{i}stico de los tests a nivel video. INCA est\'{a} bajo acuerdo institucional, no se puede compartir externamente.
  }
\end{frame}


%% ========================================
%% SLIDE 15: SCHEDULE AND FUTURE WORK
%% ========================================
% ================= FRAME DE LA PROPUESTA, FUERA EN LA DEFENSA =================
% Una defensa no tiene trabajo restante. Se deja comentado, no borrado.
% \begin{frame}{Cronograma e trabalho restante}
%   \begin{columns}[T,totalwidth=\textwidth]
%     \column{0.55\textwidth}
%       \centering
%       \footnotesize
%       \begin{tabular}{@{}llc@{}}
%         \toprule
%         \textbf{Tarefa} & \textbf{Status} & \textbf{Data} \\
%         \midrule
%         Paper submetido            & Completo & 11/06/2026 \\
%         Proposta (este documento)  & Completo & 12/06/2026 \\
%         Draft inicial da dissertação & Completo & 12/06/2026 \\
%         Análise qualitativa estendida & Pendente & 26/06/2026 \\
%         Revisões finais do manuscrito & Pendente & 26/06/2026 \\
%         Defesa final               & Pendente & 07/08/2026 \\
%         \bottomrule
%       \end{tabular}
% 
%     \column{0.45\textwidth}
%       \textbf{Trabalho restante}
%       \begin{itemize}
%         \item Análise qualitativa estendida sobre os modelos treinados
%         \item Revisões finais do manuscrito
%         \item Incorporar o feedback desta banca
%       \end{itemize}
%   \end{columns}
% 
%   \note{%
%     Estado del trabajo.
% 
%     El manuscrito del paper con la base t\'{e}cnica de la tesis ya fue submitido.
% 
%     La proposta (este documento) tambi\'{e}n est\'{a} completa.
% 
%     El draft inicial de la tesis completa est\'{a} bajo revisi\'{o}n.
% 
%     Quedan tres tareas:
%     1. An\'{a}lisis cualitativo extendido sobre los modelos ya entrenados.
%     2. Revisiones finales del manuscrito.
%     3. Incorporar el feedback que reciba de esta banca examinadora.
% 
%     El trabajo experimental ya est\'{a} terminado. Lo que falta es an\'{a}lisis y escritura.
%   }
% \end{frame}


%% ========================================
%% SLIDE 16: CONCLUSIONS
%% ========================================
\begin{frame}{Conclusões}
  \begin{itemize}
    \item SSL clássico (PL, MT) melhora a segmentação das vértebras cervicais em VFSS quando o baseline supervisionado não saturou (abaixo de F1 $\approx$ 0,85 nos regimes avaliados)
    \medskip
    \item O tamanho do pool depende do método e do dataset (em UNM, PL com pool pequeno e MT com pool completo; em INCA reduzido, esse padrão não se manteve)
    \medskip
    \item A proximidade temporal não traz vantagem sistemática sobre seleção aleatória do mesmo tamanho
  \end{itemize}

  \note{%
    Tres mensajes para llevarse de la presentaci\'{o}n.

    Primero: SSL no es universal. Funciona en r\'{e}gimen de pocos labels (UNM, INCA 10\%), no cuando la supervisada ya satur\'{o}.

    Segundo: el tama\~{n}o del pool importa, y los dos m\'{e}todos reaccionan distinto seg\'{u}n el dataset.

    Tercero: la proximidad temporal no aporta ventaja sobre random.
  }
\end{frame}


%% ========================================
%% THANK YOU (al final)
%% ========================================
\begin{frame}
  \centering
  \vfill
  {\Huge \textbf{Obrigado!}}\\[0.6em]
  {\large Perguntas?}\\[1.4em]
  \href{mailto:squispearias@aluno.puc-rio.br}{\texttt{squispearias@aluno.puc-rio.br}}\\[1.2em]
\end{frame}


%% ========================================
%% BACKUP SLIDES (no contar en el total del footer)
%% ========================================

% Guardar el total antes de los backup
\setcounter{finalframe}{\value{framenumber}}

%% ----- BACKUP 1: Training details -----
% Truco para que los backup slides NO cuenten en el total del footer
\setcounter{framenumber}{\value{finalframe}}

\begin{frame}{Backup: detalhes do treinamento}
  \begin{columns}[T,totalwidth=\textwidth]
    \column{0.5\textwidth}
      \textbf{Arquitetura}
      \begin{itemize}
        \item U-Net++ + EfficientNet-B3 (ImageNet)
        \item Entrada: $320 \times 320$, escala de cinza
        \item Loss: BCE + Dice
      \end{itemize}

      \bigskip

      \textbf{Otimização}
      \begin{itemize}
        \item Adam, lr $= 10^{-3}$, weight decay $10^{-4}$
        \item Lote: 5 rotuladas + 1 sem rótulo
        \item Early stopping em IoU de validação, paciência 40
      \end{itemize}

    \column{0.5\textwidth}
      \textbf{Detalhes do SSL}
      \begin{itemize}
        \item Teacher: EMA dos pesos do student ($\alpha = 0,99$)
        \item Threshold de PL $\tau = 0,95$ (confiança em $p$ ou $1-p$)
        \item Peso da loss não supervisionada: warmup linear
      \end{itemize}

      \bigskip

      \textbf{Augmentation}
      \begin{itemize}
        \item Fraco: flip, crop pequeno
        \item Forte: + intensity jitter, ruído gaussiano
      \end{itemize}
  \end{columns}
\end{frame}


%% ----- BACKUP 3: Other metrics -----
\begin{frame}{Backup: outras métricas (UNM full)}
  \centering
  \footnotesize
  \begin{tabular}{@{}lcc@{}}
    \toprule
    \textbf{Método}    & \textbf{Dice (F1)}    & \textbf{IoU}      \\
    \midrule
    Supervisionado     & $0,801 \pm 0,014$     & $0,702 \pm 0,010$ \\
    PL ($r{=}3$)       & $0,852 \pm 0,019$     & $0,762 \pm 0,021$ \\
    MT (pool completo) & $0,860 \pm 0,006$     & $0,769 \pm 0,007$ \\
% As colunas de Precisao e Recall foram retiradas: nao constam em nenhum
% run_report.json, portanto nao sao rastreaveis. Valores antigos, por se acaso:
%    Supervisionado  precisao 0,838  recall 0,778
%    PL (r=3)        precisao 0,866  recall 0,844
%    MT (pool)       precisao 0,871  recall 0,853
    \bottomrule
  \end{tabular}

  \vspace{1em}
  \footnotesize
  \textit{Média $\pm$ std entre seeds de treinamento. Test set: 7 vídeos, 63 frames.}
\end{frame}


%% ----- BACKUP 4: Hyperparameter sensitivity -----
\begin{frame}{Backup: divisão dos datasets}
  \begin{columns}[T,totalwidth=\textwidth]
    \column{0.5\textwidth}
      \textbf{UNM}
      \begin{itemize}
        \item Divisão por paciente (sem vazamento)
        \item Treino: 23 pacientes, 218 frames anotados
        \item Validação: 5 pacientes
        \item Teste: 7 pacientes, 63 frames
      \end{itemize}

    \column{0.5\textwidth}
      \textbf{INCA}
      \begin{itemize}
        \item Divisão por paciente
        \item Treino: 75 pacientes, 634 frames anotados
        \item Validação: 17 pacientes
        \item Teste: 23 pacientes, 194 frames
      \end{itemize}
  \end{columns}

  \vspace{1em}
  \footnotesize
  \textit{Ambos datasets seguem divisão a nível de paciente para evitar vazamento temporal entre conjuntos.}
\end{frame}


\begin{frame}{Backup: o percurso de um frame}
  \centering
  \resizebox{!}{0.72\textheight}{%
  \begin{tikzpicture}[
    stg/.style={draw=PUCBlue!55, rounded corners=3pt, fill=white,
                minimum width=3.0cm, minimum height=1.5cm,
                text width=2.8cm, align=center, font=\scriptsize, inner sep=3pt},
    net/.style={stg, fill=sslpurple, draw=Accent3!70},
    avali/.style={stg, fill=lowlabel, draw=Accent1!70},
    sslb/.style={draw=Accent2!55, rounded corners=3pt, fill=highlabel,
                 minimum width=3.7cm, minimum height=1.55cm,
                 text width=3.55cm, align=center, font=\scriptsize, inner sep=3pt,
                 execute at begin node={\hyphenpenalty=10000\exhyphenpenalty=10000}},
    fl/.style={-{Stealth[length=4pt]}, thick, PUCBlue!70},
    sf/.style={-{Stealth[length=4pt]}, thick, Accent2!75, dashed}
  ]

  % linha 1: entrada, pré-processamento, rede
  \node[stg] (n1) at (0,0)
    {\textbf{Frame lateral}\\escala de cinza};
  \node[stg, right=0.8cm of n1] (n2)
    {\textbf{Pré-processamento}\\padding, $320\times320$\\normalizado $[0,1]$\\1 canal replicado $\times 3$};
  \node[net, right=0.8cm of n2] (n3)
    {\textbf{U-Net++}\\EfficientNet-B3\\encoder ImageNet};

  \draw[fl] (n1) -- (n2);
  \draw[fl] (n2) -- (n3);

  % linha 2: saída e avaliação, da direita para a esquerda
  \node[stg, below=1.6cm of n3] (n4)
    {\textbf{Probabilidades}\\por pixel, $320\times320$};
  \node[stg, left=0.8cm of n4] (n5)
    {\textbf{Limiar 0,5}\\máscara binária};
  \node[avali, left=0.8cm of n5] (n6)
    {\textbf{F1 (Dice)}\\ASSD e HD95,\\na grade do modelo};

  \draw[fl] (n3) -- (n4);
  \draw[fl] (n4) -- (n5);
  \draw[fl] (n5) -- (n6);

  % linha 3: ramo semi-supervisionado, da esquerda para a direita
  \node[sslb, below=1.7cm of n6] (s1)
    {\textbf{Pool sem rótulo}\\mesmos vídeos de treino};
  \node[sslb, right=0.8cm of s1] (s2)
    {\textbf{Vistas fraca e forte}\\fraca: Teacher, forte: Student\\as duas passam pela rede};
  \node[sslb, right=0.8cm of s2] (s3)
    {\textbf{Perda não supervisionada}\\PL: $p\geq0{,}95$ ou $p\leq0{,}05$\\MT: consistência em todos os pixels};

  \draw[sf] (s1) -- (s2);
  \draw[sf] (s2) -- (s3);
  \draw[sf] (s3.east) -- ++(0.9,0) |- (n3.east);

  \node[draw=Accent2!45, dashed, rounded corners=4pt, inner sep=9pt,
        fit=(s1)(s2)(s3)] (panel) {};
  \node[font=\scriptsize\itshape, text=Accent2!85!black, anchor=north west]
    at ($(panel.north west) + (0.1,0.36)$) {somente durante o treino semi-supervisionado};

  \end{tikzpicture}%
  }

  \note{%
    Slide de respaldo. Solo si preguntan por el flujo, que es lo que pide siempre quien abre el trabajo por primera vez.

    La l\'{i}nea azul es lo que pasa con UN frame: entra en escala de grises, padding a casi cuadrado, 320 por 320, normalizado al intervalo unidad, y el canal se replica tres veces porque el encoder viene de ImageNet y espera tres canales. NO hay informaci\'{o}n de color: los tres canales son la misma intensidad.

    Sale un mapa de probabilidad por p\'{i}xel, se umbraliza en 0,5, y la m\'{a}scara binaria se compara con la anotaci\'{o}n por F1, que es el Dice, y por las dos m\'{e}tricas de borde, ASSD y HD95. En la rejilla del modelo, no en la resoluci\'{o}n original.

    OJO: el IoU NO es una m\'{e}trica de test en la tesis. Es el criterio de early stopping sobre validaci\'{o}n. En el deck hay IoU de test en el respaldo 29, pero esos valores no est\'{a}n en el documento.

    OJO 2: el umbral de 0,5 de la binarizaci\'{o}n no aparece en el protocolo de evaluaci\'{o}n de la tesis; el \'{u}nico 0,5 escrito est\'{a} dentro de Pseudo-Labeling, l. 738. Si preguntan de d\'{o}nde sale, es el punto medio y no se ajust\'{o}.

    La rama naranja discontinua solo existe en entrenamiento semi-supervisado: el pool sin etiquetar, las dos vistas, y la p\'{e}rdida no supervisada que entra en la misma red.

    Si preguntan por la resoluci\'{o}n de la evaluaci\'{o}n, esa es la tarjeta 17.
  }
\end{frame}

\end{document}
